\documentclass[10pt, a4paper]{article}
\usepackage{preamble}

\title{Algebra II \\
    \large Prereading Notes}
\author{Luke Phillips}
\date{October 2025}

\begin{document}

\maketitle

\newpage

\tableofcontents

\newpage

\section{Rings, subrings and fields}

\subsection{First examples
(and non-examples)
of rings and fields}

The integers $\Z = \{0, \pm 1, \pm 2, \dotsc\}$.
We can add,
subtract and multiply two integers while always remaining inside $\Z$.

Polynomials:
Let $\Q[x] = \{a_0 + a_1x + \dotsc + a_nx ^ n \mid a_i \in \Q\}$ denote the set of polynomials with rational coefficients.
We can add
subtract,
and multiply two polynomials while remaining in $\Q[x]$.

These examples lead us to an informal definition of a ring:
A ring is a set of mathematical objects
(e.g. numbers,
matrices,
functions etc.)
where we can add,
subtract and multiply two elements while remaining in the set.
If we can also divide,
the ring is called a field.

Here is a non-example of a ring
Let
\[
\N = \{1, 2, 3, \dotsc\}
\]
be the set of natural numbers.
We can add two numbers such that the result is still in $\N$.
The same holds for multiplication.
However,
we cannot subtract without eventually ending up outside of $\N$
(e.g. $1 - 3 = -2 \notin \N$).
Thus $\N$ is not a ring.


\subsection{Definition of a ring}
Now we will precisely state what is means to be a ring.

First we need to state what a binary operation is.

If we let $R \times R = \{(x, y) \mid x, y \in R\}$ be the Cartesian product,
then addition in $R$ is nothing but a function $+ : R \times R \to \R$,
but instead of $+(x, y)$ we write $x + y$.
Such functions are called binary operations,
because they take two elements as arguments.

We can recall that $R$ is an abelian group with respect to $+$ if:
\begin{enumerate}[label = (\roman*)]
    \item $R$ has an additive identity element,
    such that
    \[
    x + 0 = 0 + x = x,\quad\forall x \in R.
    \]

    \item (Associativity)
    \[
    (x + y) + z = x + (y + z),\quad \forall x, y, z \in R.
    \]

    \item $\forall x \in R, \exists y \in R$ such that
    \[
    x + y = y + x = 0.
    \]
    Such a $y$ is unique and is usually denotes as $-x$.

    \item The above ensure that $R$ is a group with respect to $+$.
    For $R$ to be abelian,
    we further require that
    \[
    x + y = y + x,\quad\forall x, y \in R.
    \]
\end{enumerate}

For the definition of a ring,
we need another binary operation,
which we usually call multiplication in $R$.
It is a function $\cdot : R \times R \to R$,
written $x \cdot y$,
or simply $xy$.

\begin{definition}[Ring]
    A ring $R$ is a set together with two binary operations called "addition" and "multiplication",
    such that $+$ makes $R$ into an abelian group,
    and such that the following conditions hold:
    \begin{enumerate}[label = \roman*)]
        \item (Identity)
        \[
        \exists 1 \in R, 1 \cdot x = x \cdot 1 = x,\quad\forall x \in R.
        \]

        \item (Associativity)
        For any $x, y, z \in R$ we have $x(yz) = (xy)z$.

        \item (Distributivity)
        For any $x, y, z \in R$ we have $x(y + z) = xy + xz$ and $(y + z)x = yx + zx$.
    \end{enumerate}
\end{definition}

If the two conditions for distributivity coincide then we call the ring commutative.

\subsection{The integers modulo \texorpdfstring{$n$}{}}
Here is an important example of rings.
Through residue classes of integers modulo some fixed natural number $n \in \N$ we can get rings.
Dividing an integer by $n$ always gets a remainder along the following numbers
\[
0, 1, \dotsc, n - 1.
\]
When an integer $a \in \Z$ is thought of as a remainder modulo $n$,
we write it as $\bar{a}$.

$\bar{a}$ is not a precise definition as $\bar{a}$ cannot be an integer,
for example when $n = 2$,
every integer has remainder $0$ or $1$,
so $\bar{0} = \bar{2} = \bar{4} = \dotsi$ and $\bar{1} = \bar{3} = \dotsi$.
But we certainly don't have $0 = 2$ so we cannot have $\bar{a} = a$.

For general $a$ and $n$,
\[
\bar{a} = \{x \in \Z \mid x \equiv a\pmod{n}\}.
\]
The elements $\bar{a}$ are called residue classes
(modulo $n$).
There are precisely $n$ of them,
\[
\{\bar{0}, \bar{1}, \dotsc, \overline{n - 1}\}
\]
we will denote this set by $\Z / n$.

\begin{example}
    We will show that $\Z / n$ is a commutative ring.
    To do this we first need to say what we mean by "adding" and "multiplying" two residue classes.
    The way to do this is by defining
    \begin{align*}
        \bar{a} + \bar{b} &= \overline{a + b}, \\
        \bar{a} \cdot \bar{b} &= \overline{ab}.
    \end{align*}
    Thus,
    to define $\bar{a} + \bar{b}$,
    we have chosen $a \in \bar{a}$ and $b \in \bar{b}$,
    taken $a + b \in \Z$ and then $\overline{a + b} \in \Z / n$.

    We will leave showing that $\Z / n$ forms an abelian group under addition since this is not too important.

    Existence of an identity element,
    commutativity,
    associativity and distributivity follow straight-forwardly,
    using the properties of multiplication in $\Z$.
    For example:
    \[
    \bar{a}(\bar{b} + \bar{c}) = \bar{a}(\overline{b + c}) = \overline{a(b + c)} = \overline{ab + ac} = \overline{ab} + \overline{ac} = \bar{a}\cdot\bar{b} + \bar{a}\cdot\bar{c},
    \]
    where we have used distributivity in the ring $\Z$ for the third equality.
\end{example}

\subsection{Subrings}

\begin{definition}[Subring]
    A subring $S$ of a ring $(R, +, \cdot)$ is a set $S \subseteq R$ such that
    \begin{enumerate}[label = \roman*)]
        \item $0 \in S, 1 \in S$,

        \item $a, b \in S \implies a + b \in S$,

        \item $a, b \in S \implies ab \in S$,

        \item $a \in S \implies -a \in S$.
    \end{enumerate}
\end{definition}

We can intuitively sum up this definition as:
"A subring contains $0$ and $1$ and is closed under addition,
subtraction and multiplication".
A subring is itself a ring.

\begin{example}
    Show that $\Z[\sqrt{2}] \coloneqq \{a + b\sqrt{2}\mid a, b \in \Z\}$ is a ring.

    \begin{solution}
        We can quickly do this by showing that $\Z[\sqrt{2}]$ is a subring of $\R$.
        $0, 1 \in \Z[\sqrt{2}]$ can be seen easily.
        \begin{align*}
            (a + b\sqrt{2}) + (c + d\sqrt{2}) &= (a + c) + (b + d)\sqrt{2} \in \Z[\sqrt{2}], \\
            (a + b\sqrt{2})(c + d\sqrt{2}) &= (ac + 2bd) + (ab + bc)\sqrt{2} \in \Z[\sqrt{2}], \\
            -(a + b\sqrt{2}) &= -a - b\sqrt{2} \in \Z[\sqrt{2}].
        \end{align*}
        Therefore $\Z[\sqrt{2}]$ is a subring of $\R$ and is hence a ring.
    \end{solution}
\end{example}

\subsection{Fields}

\begin{definition}[Fields]
    A ring $(R, +, \cdot)$ is called a field if
    \begin{enumerate}[label = \roman*)]
        \item $R$ is commutative,

        \item $1 \neq 0$ in $R$.

        \item For any $a \in R$ such that $a \neq 0$,
        there exists a $b \in R$ such that $ab = 1$
        ($b$ is called the inverse of $a$ and we usually write $b = a ^ {-1}$).
    \end{enumerate}
\end{definition}
Inverses are unique:
If $ab = 1$ and $ab' = 1$,
then $abb' = b'$,
so
\[
b' = abb' = bab' = b \cdot 1 = b.
\]

\newpage

\section{Arithmetic in \texorpdfstring{$\Z$}{} and \texorpdfstring{$F[x]$}{}}
In a commutative ring $R$,
an element $a \in R$ is said to divide an element $b \in R$,
written $a \mid b$,
if there exists an $x \in R$ such that
\[
b = ax.
\]

\subsection{Divisibility in \texorpdfstring{$\Z$}{}}

\begin{lemma}[Well-Ordering Principle
    (WOP)]
    Let $S$ be a non-empty subset of $\N_0 = \N \cup \{0\}$.
    Then $S$ has a smallest element.
\end{lemma}

\textbf{Division with remainder in $\Z$}

\begin{example}
    Find the integer quotient and remainder of $435$ divided by $43$ in $\Z$

    \begin{solution}
        \[
        435 = 43 \cdot 10 + 5,
        \]
        and $5 < 43$ so we stop.
        The quotient is $q = 10$ and remainder $r = 5$.
    \end{solution}
\end{example}

\begin{proposition}
    Let $a \in \Z$,
    $b \neq 0$.
    Then there exist integers $q$ and $r$ such that
    \[
    a = q \cdot b + r,\quad\text{and}\quad 0 \leq r < |b|.
    \]

    \begin{proof}
        First,
        we can assume that $b > 0$.
        Indeed,
        if $b < 0$,
        then we can apply the statement to $-b > 0$ and use $-q$ instead of $q$.
        This does not change $|b|$ so the inequalities still hold.
        Consider the subsets
        \[
        S ^ {+} = \{a + bk\mid k \in \Z, a + bk \geq 0\} \subset S = \{a + bk \mid k \in \Z\}.
        \]
        $S ^ {+}$ is non-empty since both
        \[
        a + b \cdot 0 = a \in S\quad\text{and}\quad a + b(-a) \in S
        \]
        and one of these is always $\geq 0$
        (if $a < 0$,
        then $a + b(-a) = (-a)(b - 1) \geq 0$).

        So by the well ordering principle,
        $S ^ {+}$ contains a minimal element $r = a + bk$,
        for some $k$.
        Also,
        \[
        r - b = a + b(k - 1) \in S,
        \]
        but this cannot lie in $S ^ {+}$ because $r - b < r$ and $r$ is minimal in $S ^ {+}$.
        Thus $r - b < 0$,
        so $r < b$.
    \end{proof}
\end{proposition}

\textbf{The Euclidean Algorithm in $\Z$}

\begin{theorem}
    Let $a, b \in \Z$,
    not both zero.
    Then $\gcd(a, b)$ exists and can be computed by the Euclidean algorithm.
    Moreover,
    if $d = \gcd(a, b)$,
    then these exist $x, y \in \Z$ such that
    \[
    ax + by = d.
    \]
\end{theorem}

\begin{example}
    Use the Euclidean Algorithm to compute $d = \gcd(455, 1235)$ and find integers $x, y$ such that $d = 455x + 1235y$.

    \begin{solution}
        Division with remainder:
        \begin{align*}
            1235 &= 2 \cdot 455 + 325 \\
            455 &= 1 \cdot 325 + 130 \\325 &= 2 \cdot 130 + 65 \\
            130 &= 2 \cdot 65.
        \end{align*}
        When the remainder is $0$ we stop.
        The $\gcd$ will be the last non-zero remainder,
        here $d = 65$.

        Solving for the remainders,
        substituting back and simplifying,
        we get
        \begin{align*}
            65 &= 325 - 2 \cdot 130 \\
            &= 325 - 2(455 - 1 \cdot 325) \\
            &= \dotsi \\
            &= 455 \cdot (-8) + 1235 \cdot 3
        \end{align*}
        so $x = -8$ and $y = 3$.
    \end{solution}
\end{example}

\textbf{Primes and unique factorisation in $\Z$}

\begin{lemma}
    If $p$ is a prime and $p \mid ab$ for some $a, b \in \Z$,
    then $p \mid a$ or $p \mid b$
    (in other words $p$ prime implies $(\forall a, b \in \Z, p \mid ab \implies p \mid a\  \wedge\ p \mid b$)
\end{lemma}

\begin{theorem}[Fundamental Theorem of Arithmetic
    (FTA)]
    Let $n \in \N$,
    $n > 1$.
    Then $n$ can be written as a product of primes in a unique way.
    More precisely,
    there are distinct primes $p_1, \dotsc, p_r$ and $e_1, \dotsc, e_r \in \N$ such that
    \[
    n = p_1 ^ {e_1} \dotsi p_r ^ {e_r},
    \]
    and if $n = q_1 ^ {f_1} \dotsi q_s ^ {f_s}$ for some distinct primes $q_1, \dotsc, q_s$ and $f_1, \dotsc, f_s \in \N$,
    then $r = s$,
    and up to re-ordering,
    $p_i = q_i$ and $e_i = f_i$ for all $i$.
    
\end{theorem}




\end{document}